% Figure: the five series resistances of the kerosene-cooler overall coefficient,
% each referred to the tube outer area, drawn as a stacked bar so the reader can
% read at a glance which term governs. The kerosene-side film and its fouling
% together are the tall pair; the water-side film is a sliver.
\begin{figure}[htbp]
\centering
\begin{tikzpicture}
\begin{axis}[
  width=12cm, height=6.5cm,
  xbar stacked,
  xmin=0, xmax=2.6,
  ytick={1},
  yticklabels={},
  y=1.4cm,
  xlabel={resistance referred to tube outer area ($10^{-3}\,\text{m}^2\cdot\text{K}/\text{W}$)},
  tick label style={font=\scriptsize},
  label style={font=\small},
  legend style={font=\scriptsize, at={(0.5,-0.28)}, anchor=north, legend columns=2},
  legend cell align=left,
  enlarge y limits=0.6,
]
% values in 1e-3 m2K/W: shell film 0.200, shell fouling 0.300, wall 0.045,
% tube fouling referred out 0.507, tube film referred out 1.407
\addplot[fill=engblue] coordinates {(0.200,1)};
\addplot[fill=engamber] coordinates {(0.300,1)};
\addplot[fill=enggray] coordinates {(0.045,1)};
\addplot[fill=engorange] coordinates {(0.507,1)};
\addplot[fill=engred] coordinates {(1.407,1)};
\legend{water film, water fouling, tube wall, kerosene fouling, kerosene film}
\end{axis}
\end{tikzpicture}
\caption[Term-by-term resistance stack of the kerosene cooler]{The five series
resistances of the kerosene cooler's overall coefficient, each referred to the tube
outer area and drawn to scale. The kerosene-side film and the kerosene-side fouling
together fill the right two-thirds of the bar, so the tube side governs and any
money spent to improve the exchanger belongs there. The water-side film, despite a
respectable film coefficient, is the leftmost sliver because water carries heat far
better than the oil it works against. Reading the stack rather than a single
tabulated overall coefficient is what tells the designer where the bottleneck sits
and therefore which change would move it.}
\label{fig:vol04ch06:resistance-stack}
\end{figure}
